\nwfilename{nw/many-sorted-term.nw}\nwbegindocs{0}% -*- mode: poly-noweb; noweb-code-mode: sml-mode; -*-% ===> this file was generated automatically by noweave --- better not edit it
%%
%% many-sorted-term.nw
%% 
%% Made by Alex Nelson <pqnelson@gmail.com>
%% Login   <alex@lisp>
%% 
%% Started on  2026-04-04T18:09:28-0700
%% Last update 2026-04-04T18:09:28-0700
%% 

\section{Syntax for terms}

\begin{node}
Since a proof assistant resembles an interpreter in many ways, it
should not be surprising that we will first focus on implement
abstract syntax trees for terms and formulas for three-sorted
first-order logic. 
\end{node}

\begin{node}
The grammar for a term in $FS_{0}$'s three-sorted logic is, more or
less, the same as in single-sorted logic. Since we are working with
\emph{three} sorts, we isolate it in is own module. We would have the
BNF for such things be:
\begin{center}
\begin{tabular}{rcl}
$\langle$\textit{sort}$\rangle$ & $::=$ & \texttt{Ind} $\mid$ \texttt{Fun} $\mid$ \texttt{Class}\\
$\langle$\textit{term}$\rangle$ & $::=$ & $\langle$\textit{variable}$\rangle$ \\
& $\mid$ & $\langle$\textit{function}$\rangle$
\end{tabular}
\end{center}
Remember, a ``function'' here should be thought of as a ``term constructor''.

We've abstracted away the ``concrete syntax'' for the sorted
first-order logic, so we can treat variables as an ordered pair
consisting of a lexeme (string representation of its identifier) and
its sort. Similarly, we can treat a function as an ordered triple
consisting of a lexeme, its sort, and a list of terms for its
arguments.

Before getting too far ahead of ourselves, let us implement a simple
encoding for the sort. We will just make everything public, because
there's not much to hide. There are three data constructors (one for
each sort), a comparison function, and a function to get the string
representation for the sort.

Also we should bear in mind the user may have a ``mismatch'' between
the sorts expected for arguments of predicates and functions, and the
arguments supplied for them. This situation should raise an
exception. We will create a {\Tt{}Sort.Mismatch\nwendquote} exception, with a string
message for what went wrong.
\end{node}

\nwenddocs{}\nwbegincode{1}\sublabel{NW1rTmOJ-3ZqggC-1}\nwmargintag{{\nwtagstyle{}\subpageref{NW1rTmOJ-3ZqggC-1}}}\moddef{Sort.sml~{\nwtagstyle{}\subpageref{NW1rTmOJ-3ZqggC-1}}}\endmoddef\nwstartdeflinemarkup\nwenddeflinemarkup
structure Sort = struct
exception Mismatch of string;
datatype t = IND | FUN | CLASS;

(* compare : \nwlinkedidentc{Sort.t}{NW1rTmOJ-3ZqggC-1} * \nwlinkedidentc{Sort.t}{NW1rTmOJ-3ZqggC-1} -> order *)
fun compare(IND,IND) = EQUAL
  | compare(IND, _) = LESS
  | compare(FUN,IND) = GREATER
  | compare(FUN,FUN) = EQUAL
  | compare(FUN, _) = LESS
  | compare(CLASS, CLASS) = EQUAL
  | compare(CLASS, _) = GREATER;

(* to_string : \nwlinkedidentc{Sort.t}{NW1rTmOJ-3ZqggC-1} -> string *)
fun to_string (IND) = "IND"
  | to_string (FUN) = "FUN"
  | to_string (CLASS) = "CLASS";
end;
\nwindexdefn{\nwixident{Sort.t}}{Sort.t}{NW1rTmOJ-3ZqggC-1}\nwindexdefn{\nwixident{Sort.compare}}{Sort.compare}{NW1rTmOJ-3ZqggC-1}\nwindexdefn{\nwixident{Sort.to{\_}string}}{Sort.to:unstring}{NW1rTmOJ-3ZqggC-1}\nwindexdefn{\nwixident{Sort.IND}}{Sort.IND}{NW1rTmOJ-3ZqggC-1}\nwindexdefn{\nwixident{Sort.FUN}}{Sort.FUN}{NW1rTmOJ-3ZqggC-1}\nwindexdefn{\nwixident{Sort.CLASS}}{Sort.CLASS}{NW1rTmOJ-3ZqggC-1}\eatline
\nwnotused{Sort.sml}\nwidentdefs{\\{{\nwixident{Sort.CLASS}}{Sort.CLASS}}\\{{\nwixident{Sort.compare}}{Sort.compare}}\\{{\nwixident{Sort.FUN}}{Sort.FUN}}\\{{\nwixident{Sort.IND}}{Sort.IND}}\\{{\nwixident{Sort.t}}{Sort.t}}\\{{\nwixident{Sort.to{\_}string}}{Sort.to:unstring}}}\nwendcode{}\nwbegindocs{2}\nwdocspar
\begin{node}
Now, a three-sorted term is either a variable (which is a string and a
sort) or a function (which is a string, a sort, and a list of arguments).
We do not ``hardcode'' any of the primitive notions from $FS_{0}$ in
this representation. Instead, they will be ``initialized'' in the
state of the kernel.
\end{node}

\nwenddocs{}\nwbegincode{3}\sublabel{NW1rTmOJ-4LsRWC-1}\nwmargintag{{\nwtagstyle{}\subpageref{NW1rTmOJ-4LsRWC-1}}}\moddef{Term.sml~{\nwtagstyle{}\subpageref{NW1rTmOJ-4LsRWC-1}}}\endmoddef\nwstartdeflinemarkup\nwenddeflinemarkup
structure Term = struct
datatype t = Var of string * \nwlinkedidentc{Sort.t}{NW1rTmOJ-3ZqggC-1}
           | Fun of string * \nwlinkedidentc{Sort.t}{NW1rTmOJ-3ZqggC-1} * t list;

\LA{}Test if two terms have the same sort~{\nwtagstyle{}\subpageref{NW1rTmOJ-1WQNps-1}}\RA{}

\LA{}Predicates of terms about their sorts~{\nwtagstyle{}\subpageref{NW1rTmOJ-3FfXTO-1}}\RA{}

\LA{}Free variables in a term~{\nwtagstyle{}\subpageref{NW1rTmOJ-2nRtky-1}}\RA{}

\LA{}Substitution of a term for a variable in a term~{\nwtagstyle{}\subpageref{NW1rTmOJ-3OIyM5-1}}\RA{}

\LA{}Comparing two terms~{\nwtagstyle{}\subpageref{NW1rTmOJ-2yJYsC-1}}\RA{}

\LA{}Equality of two terms~{\nwtagstyle{}\subpageref{NW1rTmOJ-3810sY-1}}\RA{}

\LA{}Test if term contains subterm~{\nwtagstyle{}\subpageref{NW1rTmOJ-Re21r-1}}\RA{}

\LA{}Determine a fresh variable given some free variables~{\nwtagstyle{}\subpageref{NW1rTmOJ-4Ykq9M-1}}\RA{}

\LA{}Serialize a term~{\nwtagstyle{}\subpageref{NW1rTmOJ-40h2U3-1}}\RA{}
end;
\nwindexdefn{\nwixident{Term.t}}{Term.t}{NW1rTmOJ-4LsRWC-1}\eatline
\nwnotused{Term.sml}\nwidentdefs{\\{{\nwixident{Term.t}}{Term.t}}}\nwidentuses{\\{{\nwixident{Sort.t}}{Sort.t}}}\nwindexuse{\nwixident{Sort.t}}{Sort.t}{NW1rTmOJ-4LsRWC-1}\nwendcode{}\nwbegindocs{4}\nwdocspar
\begin{node}
We provide a minimal abstraction for terms. As it stands now, without
any ``helper'' or ``uniform'' constructors, it is about 44 lines of
code. If we ``hardcoded'' the $FS_{0}$ primitives, that would require
about 412 lines of code --- about $9.36$ times as much code.
\end{node}

\begin{node}
We provide a uniform way to access the sort of a term, and test if two
terms have the same sort.
\end{node}

\nwenddocs{}\nwbegincode{5}\sublabel{NW1rTmOJ-1WQNps-1}\nwmargintag{{\nwtagstyle{}\subpageref{NW1rTmOJ-1WQNps-1}}}\moddef{Test if two terms have the same sort~{\nwtagstyle{}\subpageref{NW1rTmOJ-1WQNps-1}}}\endmoddef\nwstartdeflinemarkup\nwusesondefline{\\{NW1rTmOJ-4LsRWC-1}}\nwenddeflinemarkup
(* sort : \nwlinkedidentc{Term.t}{NW1rTmOJ-4LsRWC-1} -> \nwlinkedidentc{Sort.t}{NW1rTmOJ-3ZqggC-1} *)
fun sort (Var (_, s)) = s
  | sort (Fun (_, s, _)) = s;

(* have_same_sort : \nwlinkedidentc{Term.t}{NW1rTmOJ-4LsRWC-1} -> \nwlinkedidentc{Term.t}{NW1rTmOJ-4LsRWC-1} -> bool *)
fun have_same_sort (lhs : t) (rhs : t) =
  (sort lhs) = (sort rhs);
\nwindexdefn{\nwixident{Term.sort}}{Term.sort}{NW1rTmOJ-1WQNps-1}\nwindexdefn{\nwixident{Term.have{\_}same{\_}sort}}{Term.have:unsame:unsort}{NW1rTmOJ-1WQNps-1}\eatline
\nwused{\\{NW1rTmOJ-4LsRWC-1}}\nwidentdefs{\\{{\nwixident{Term.have{\_}same{\_}sort}}{Term.have:unsame:unsort}}\\{{\nwixident{Term.sort}}{Term.sort}}}\nwidentuses{\\{{\nwixident{Sort.t}}{Sort.t}}\\{{\nwixident{Term.t}}{Term.t}}}\nwindexuse{\nwixident{Sort.t}}{Sort.t}{NW1rTmOJ-1WQNps-1}\nwindexuse{\nwixident{Term.t}}{Term.t}{NW1rTmOJ-1WQNps-1}\nwendcode{}\nwbegindocs{6}\nwdocspar
\begin{node}
We have some helper functions for common predicates.
\end{node}

\nwenddocs{}\nwbegincode{7}\sublabel{NW1rTmOJ-3FfXTO-1}\nwmargintag{{\nwtagstyle{}\subpageref{NW1rTmOJ-3FfXTO-1}}}\moddef{Predicates of terms about their sorts~{\nwtagstyle{}\subpageref{NW1rTmOJ-3FfXTO-1}}}\endmoddef\nwstartdeflinemarkup\nwusesondefline{\\{NW1rTmOJ-4LsRWC-1}}\nwenddeflinemarkup
(* is_var : \nwlinkedidentc{Term.t}{NW1rTmOJ-4LsRWC-1} -> bool *)
fun is_var (Var _) = true
  | is_var _ = false;

(* is_ind : \nwlinkedidentc{Term.t}{NW1rTmOJ-4LsRWC-1} -> bool *)
fun is_ind (tm : t) = (\nwlinkedidentc{Sort.IND}{NW1rTmOJ-3ZqggC-1} = sort tm);

(* is_fun : \nwlinkedidentc{Term.t}{NW1rTmOJ-4LsRWC-1} -> bool *)
fun is_fun (tm : t) = (\nwlinkedidentc{Sort.FUN}{NW1rTmOJ-3ZqggC-1} = sort tm);

(* is_class : \nwlinkedidentc{Term.t}{NW1rTmOJ-4LsRWC-1} -> bool *)
fun is_class (tm : t) = (\nwlinkedidentc{Sort.CLASS}{NW1rTmOJ-3ZqggC-1} = sort tm);
\nwindexdefn{\nwixident{Term.is{\_}ind}}{Term.is:unind}{NW1rTmOJ-3FfXTO-1}\nwindexdefn{\nwixident{Term.is{\_}fun}}{Term.is:unfun}{NW1rTmOJ-3FfXTO-1}\nwindexdefn{\nwixident{Term.is{\_}class}}{Term.is:unclass}{NW1rTmOJ-3FfXTO-1}\eatline
\nwused{\\{NW1rTmOJ-4LsRWC-1}}\nwidentdefs{\\{{\nwixident{Term.is{\_}class}}{Term.is:unclass}}\\{{\nwixident{Term.is{\_}fun}}{Term.is:unfun}}\\{{\nwixident{Term.is{\_}ind}}{Term.is:unind}}}\nwidentuses{\\{{\nwixident{Sort.CLASS}}{Sort.CLASS}}\\{{\nwixident{Sort.FUN}}{Sort.FUN}}\\{{\nwixident{Sort.IND}}{Sort.IND}}\\{{\nwixident{Term.t}}{Term.t}}}\nwindexuse{\nwixident{Sort.CLASS}}{Sort.CLASS}{NW1rTmOJ-3FfXTO-1}\nwindexuse{\nwixident{Sort.FUN}}{Sort.FUN}{NW1rTmOJ-3FfXTO-1}\nwindexuse{\nwixident{Sort.IND}}{Sort.IND}{NW1rTmOJ-3FfXTO-1}\nwindexuse{\nwixident{Term.t}}{Term.t}{NW1rTmOJ-3FfXTO-1}\nwendcode{}\nwbegindocs{8}\nwdocspar
\begin{node}
The free variables of a term are just \emph{all} the variables
appearing in the term. Later, formulas will remove the bound variables.
\end{node}

\nwenddocs{}\nwbegincode{9}\sublabel{NW1rTmOJ-2nRtky-1}\nwmargintag{{\nwtagstyle{}\subpageref{NW1rTmOJ-2nRtky-1}}}\moddef{Free variables in a term~{\nwtagstyle{}\subpageref{NW1rTmOJ-2nRtky-1}}}\endmoddef\nwstartdeflinemarkup\nwusesondefline{\\{NW1rTmOJ-4LsRWC-1}}\nwenddeflinemarkup
(* fv : \nwlinkedidentc{Term.t}{NW1rTmOJ-4LsRWC-1} -> \nwlinkedidentc{Term.t}{NW1rTmOJ-4LsRWC-1} list *)
fun fv (x as (Var _)) = [x]
  | fv (Fun (_, _, args)) = List.foldr (fn (arg, fvs) => (fv arg) @ fvs)
                                       []
                                       args;
\nwindexdefn{\nwixident{Term.fv}}{Term.fv}{NW1rTmOJ-2nRtky-1}\eatline
\nwused{\\{NW1rTmOJ-4LsRWC-1}}\nwidentdefs{\\{{\nwixident{Term.fv}}{Term.fv}}}\nwidentuses{\\{{\nwixident{Term.t}}{Term.t}}}\nwindexuse{\nwixident{Term.t}}{Term.t}{NW1rTmOJ-2nRtky-1}\nwendcode{}\nwbegindocs{10}\nwdocspar
\begin{node}
Substitution amounts to replacing a variable \emph{of a given sort}
with a term \emph{of the same sort} in a term recursively by:
\begin{enumerate}
\item $x[t/x]=t$ 
\item $y[t/x]=y$ if $y\neq x$
\item $(f(t_{1},\dots,t_{n}))[t/x]=(f(t_{1}[t/x],\dots,t_{n}[t/x]))$
\end{enumerate}
Although we \emph{should} restrict our focus to the situation where
$x$ is a variable, I suspect there might be situations where we want
to replace \emph{entire subterms}.
\end{node}

\nwenddocs{}\nwbegincode{11}\sublabel{NW1rTmOJ-3OIyM5-1}\nwmargintag{{\nwtagstyle{}\subpageref{NW1rTmOJ-3OIyM5-1}}}\moddef{Substitution of a term for a variable in a term~{\nwtagstyle{}\subpageref{NW1rTmOJ-3OIyM5-1}}}\endmoddef\nwstartdeflinemarkup\nwusesondefline{\\{NW1rTmOJ-4LsRWC-1}}\nwenddeflinemarkup
(* subst : \nwlinkedidentc{Term.t}{NW1rTmOJ-4LsRWC-1} -> \nwlinkedidentc{Term.t}{NW1rTmOJ-4LsRWC-1} -> \nwlinkedidentc{Term.t}{NW1rTmOJ-4LsRWC-1} -> \nwlinkedidentc{Term.t}{NW1rTmOJ-4LsRWC-1} *)
fun subst (x : t) (replacement : t) (tm : t) =
  if not (have_same_sort x replacement)
  then raise Sort.Mismatch "\nwlinkedidentc{Term.subst}{NW1rTmOJ-3OIyM5-1}: variable and replacement have different sorts"
  else if x = tm then replacement
  else (case tm of
         (Var _) => if x = tm then replacement else tm
       | (Fun(f,s,args)) => Fun(f,s,map (subst x replacement) args));
\nwindexdefn{\nwixident{Term.subst}}{Term.subst}{NW1rTmOJ-3OIyM5-1}\eatline
\nwused{\\{NW1rTmOJ-4LsRWC-1}}\nwidentdefs{\\{{\nwixident{Term.subst}}{Term.subst}}}\nwidentuses{\\{{\nwixident{Term.t}}{Term.t}}}\nwindexuse{\nwixident{Term.t}}{Term.t}{NW1rTmOJ-3OIyM5-1}\nwendcode{}\nwbegindocs{12}\nwdocspar
\begin{node}
Comparing two terms follows the following rules:
\begin{enumerate}
\item First compare the sort of both sides; if they differ, then this
  is the result
\item Variables are less than or equal to everything
\item Functions are greater than or equal to everything
\item Variables are compared by sort, then by name
\item Functions are compared by sort, then by name, then
  lexicographically on the arguments
\end{enumerate}
\end{node}


\nwenddocs{}\nwbegincode{13}\sublabel{NW1rTmOJ-2yJYsC-1}\nwmargintag{{\nwtagstyle{}\subpageref{NW1rTmOJ-2yJYsC-1}}}\moddef{Comparing two terms~{\nwtagstyle{}\subpageref{NW1rTmOJ-2yJYsC-1}}}\endmoddef\nwstartdeflinemarkup\nwusesondefline{\\{NW1rTmOJ-4LsRWC-1}}\nwenddeflinemarkup
(* compare : \nwlinkedidentc{Term.t}{NW1rTmOJ-4LsRWC-1} * \nwlinkedidentc{Term.t}{NW1rTmOJ-4LsRWC-1} -> order *)
fun compare(Var(x1,s1), Var(x2,s2)) =
    (case \nwlinkedidentc{Sort.compare}{NW1rTmOJ-3ZqggC-1}(s1, s2) of
      EQUAL => String.compare(x1, x2)
    | r => r)
  | compare(Var _, _) = LESS
  | compare(Fun _, Var _) = GREATER
  | compare(Fun(f1,s1,args1), Fun(f2,s2,args2)) =
    (case \nwlinkedidentc{Sort.compare}{NW1rTmOJ-3ZqggC-1}(s1,s2) of
      EQUAL => (case String.compare(f1, f2) of
                 EQUAL => compare_lists args1 args2
               | r => r)
    | r => r)
(* compare_lists : \nwlinkedidentc{Term.t}{NW1rTmOJ-4LsRWC-1} list -> \nwlinkedidentc{Term.t}{NW1rTmOJ-4LsRWC-1} list -> order *)
and compare_lists [] [] = EQUAL
  | compare_lists [] ys = LESS
  | compare_lists xs [] = GREATER
  | compare_lists (x::xs) (y::ys) =
     (case compare(x,y) of
       EQUAL => compare_lists xs ys
     | r => r);
\nwindexdefn{\nwixident{Term.compare}}{Term.compare}{NW1rTmOJ-2yJYsC-1}\nwindexdefn{\nwixident{Term.compare{\_}lists}}{Term.compare:unlists}{NW1rTmOJ-2yJYsC-1}\eatline
\nwused{\\{NW1rTmOJ-4LsRWC-1}}\nwidentdefs{\\{{\nwixident{Term.compare}}{Term.compare}}\\{{\nwixident{Term.compare{\_}lists}}{Term.compare:unlists}}}\nwidentuses{\\{{\nwixident{Sort.compare}}{Sort.compare}}\\{{\nwixident{Term.t}}{Term.t}}}\nwindexuse{\nwixident{Sort.compare}}{Sort.compare}{NW1rTmOJ-2yJYsC-1}\nwindexuse{\nwixident{Term.t}}{Term.t}{NW1rTmOJ-2yJYsC-1}\nwendcode{}\nwbegindocs{14}\nwdocspar
\begin{node}
The \emph{syntactic} equality of terms can be determined in a more
performant manner. Two variables are equal if their payload are equal
(i.e., they have the same sort and the same lexeme). Two functions are
equal if they have the same name, same sort, and their arguments are
syntactically equal (determined by recursively invoking {\Tt{}\nwlinkedidentq{Term.eq}{NW1rTmOJ-3810sY-1}\nwendquote}
pairwise on the arguments).
\end{node}

\nwenddocs{}\nwbegincode{15}\sublabel{NW1rTmOJ-3810sY-1}\nwmargintag{{\nwtagstyle{}\subpageref{NW1rTmOJ-3810sY-1}}}\moddef{Equality of two terms~{\nwtagstyle{}\subpageref{NW1rTmOJ-3810sY-1}}}\endmoddef\nwstartdeflinemarkup\nwusesondefline{\\{NW1rTmOJ-4LsRWC-1}}\nwenddeflinemarkup
(* eq : \nwlinkedidentc{Term.t}{NW1rTmOJ-4LsRWC-1} -> \nwlinkedidentc{Term.t}{NW1rTmOJ-4LsRWC-1} -> bool *)
fun eq (Var(x1,s1)) (Var(x2,s2)) = (s1 = s2) andalso (x1 = x2)
  | eq (Fun(f1,s1,args1)) (Fun(f2,s2,args2)) =
      (s1 = s2) andalso (f1 = f2) andalso (eq_lists args1 args2)
  | eq _ _ = false
(* eq_lists : \nwlinkedidentc{Term.t}{NW1rTmOJ-4LsRWC-1} list -> \nwlinkedidentc{Term.t}{NW1rTmOJ-4LsRWC-1} list -> bool *)
and eq_lists [] [] = true
  | eq_lists (x::xs) (y::ys) = (eq x y) andalso (eq_lists xs ys)
  | eq_lists _ _ = false;
\nwindexdefn{\nwixident{Term.eq}}{Term.eq}{NW1rTmOJ-3810sY-1}\nwindexdefn{\nwixident{Term.eq{\_}lists}}{Term.eq:unlists}{NW1rTmOJ-3810sY-1}\eatline
\nwused{\\{NW1rTmOJ-4LsRWC-1}}\nwidentdefs{\\{{\nwixident{Term.eq}}{Term.eq}}\\{{\nwixident{Term.eq{\_}lists}}{Term.eq:unlists}}}\nwidentuses{\\{{\nwixident{Term.t}}{Term.t}}}\nwindexuse{\nwixident{Term.t}}{Term.t}{NW1rTmOJ-3810sY-1}\nwendcode{}\nwbegindocs{16}\nwdocspar
\begin{node}
We also want to check if a term contains a subterm. This is not a
``strict'' comparison: any term occurs in itself.
Our intuition should be thinking of $x\leq y$ rather than $x<y$. 
\end{node}

\nwenddocs{}\nwbegincode{17}\sublabel{NW1rTmOJ-Re21r-1}\nwmargintag{{\nwtagstyle{}\subpageref{NW1rTmOJ-Re21r-1}}}\moddef{Test if term contains subterm~{\nwtagstyle{}\subpageref{NW1rTmOJ-Re21r-1}}}\endmoddef\nwstartdeflinemarkup\nwusesondefline{\\{NW1rTmOJ-4LsRWC-1}}\nwenddeflinemarkup
(* occurs_in : \nwlinkedidentc{Term.t}{NW1rTmOJ-4LsRWC-1} -> \nwlinkedidentc{Term.t}{NW1rTmOJ-4LsRWC-1} -> bool *)
fun occurs_in (subtm : t) (tm : t) : bool =
  (eq subtm tm) orelse
  (case tm of
    (Var _) => false
  | (Fun(_,_,args)) => List.exists (occurs_in subtm) args);
\nwindexdefn{\nwixident{Term.occurs{\_}in}}{Term.occurs:unin}{NW1rTmOJ-Re21r-1}\eatline
\nwused{\\{NW1rTmOJ-4LsRWC-1}}\nwidentdefs{\\{{\nwixident{Term.occurs{\_}in}}{Term.occurs:unin}}}\nwidentuses{\\{{\nwixident{Term.t}}{Term.t}}}\nwindexuse{\nwixident{Term.t}}{Term.t}{NW1rTmOJ-Re21r-1}\nwendcode{}\nwbegindocs{18}\nwdocspar
\begin{node}
We will want to find a fresh variable which does not appear in a list
of variables.
\end{node}

\nwenddocs{}\nwbegincode{19}\sublabel{NW1rTmOJ-4Ykq9M-1}\nwmargintag{{\nwtagstyle{}\subpageref{NW1rTmOJ-4Ykq9M-1}}}\moddef{Determine a fresh variable given some free variables~{\nwtagstyle{}\subpageref{NW1rTmOJ-4Ykq9M-1}}}\endmoddef\nwstartdeflinemarkup\nwusesondefline{\\{NW1rTmOJ-4LsRWC-1}}\nwenddeflinemarkup
local
  fun fresh_var_iter (x : string) (s : \nwlinkedidentc{Sort.t}{NW1rTmOJ-3ZqggC-1}) (vars : t list) n = 
    let val var = Var(x ^ (Int.toString n), s)
    in if List.exists (eq var) vars
       then fresh_var_iter x s vars (n + 1)
       else var
    end;
in

fun fresh_var (x : string) (s : \nwlinkedidentc{Sort.t}{NW1rTmOJ-3ZqggC-1}) (vars : t list) = 
  let val var = Var(x, s)
  in if List.exists (eq var) vars
     then fresh_var_iter x s vars 0
     else var
  end;

end;
\nwindexdefn{\nwixident{Term.fresh{\_}var}}{Term.fresh:unvar}{NW1rTmOJ-4Ykq9M-1}\eatline
\nwused{\\{NW1rTmOJ-4LsRWC-1}}\nwidentdefs{\\{{\nwixident{Term.fresh{\_}var}}{Term.fresh:unvar}}}\nwidentuses{\\{{\nwixident{Sort.t}}{Sort.t}}}\nwindexuse{\nwixident{Sort.t}}{Sort.t}{NW1rTmOJ-4Ykq9M-1}\nwendcode{}\nwbegindocs{20}\nwdocspar
\begin{node}
For debugging purposes, it will be useful to ``serialize'' a term (to
produce a string representation of the code needed to construct it).
\end{node}

\nwenddocs{}\nwbegincode{21}\sublabel{NW1rTmOJ-40h2U3-1}\nwmargintag{{\nwtagstyle{}\subpageref{NW1rTmOJ-40h2U3-1}}}\moddef{Serialize a term~{\nwtagstyle{}\subpageref{NW1rTmOJ-40h2U3-1}}}\endmoddef\nwstartdeflinemarkup\nwusesondefline{\\{NW1rTmOJ-4LsRWC-1}}\nwenddeflinemarkup
fun serialize (Var(x,s)) = "Var("^x^", "^(\nwlinkedidentc{Sort.to_string}{NW1rTmOJ-3ZqggC-1} s)^")"
  | serialize (Fun(f, s, args)) = "Fun("^f^
                                  ", "^
                                  (\nwlinkedidentc{Sort.to_string}{NW1rTmOJ-3ZqggC-1} s)^
                                  ", ["^
                                  (String.concatWith ", " (map serialize args))^
                                  "])";
\nwindexdefn{\nwixident{Term.serialize}}{Term.serialize}{NW1rTmOJ-40h2U3-1}\eatline

\nwused{\\{NW1rTmOJ-4LsRWC-1}}\nwidentdefs{\\{{\nwixident{Term.serialize}}{Term.serialize}}}\nwidentuses{\\{{\nwixident{Sort.to{\_}string}}{Sort.to:unstring}}}\nwindexuse{\nwixident{Sort.to{\_}string}}{Sort.to:unstring}{NW1rTmOJ-40h2U3-1}\nwendcode{}

\nwixlogsorted{c}{{Comparing two terms}{NW1rTmOJ-2yJYsC-1}{\nwixu{NW1rTmOJ-4LsRWC-1}\nwixd{NW1rTmOJ-2yJYsC-1}}}%
\nwixlogsorted{c}{{Determine a fresh variable given some free variables}{NW1rTmOJ-4Ykq9M-1}{\nwixu{NW1rTmOJ-4LsRWC-1}\nwixd{NW1rTmOJ-4Ykq9M-1}}}%
\nwixlogsorted{c}{{Equality of two terms}{NW1rTmOJ-3810sY-1}{\nwixu{NW1rTmOJ-4LsRWC-1}\nwixd{NW1rTmOJ-3810sY-1}}}%
\nwixlogsorted{c}{{Free variables in a term}{NW1rTmOJ-2nRtky-1}{\nwixu{NW1rTmOJ-4LsRWC-1}\nwixd{NW1rTmOJ-2nRtky-1}}}%
\nwixlogsorted{c}{{Predicates of terms about their sorts}{NW1rTmOJ-3FfXTO-1}{\nwixu{NW1rTmOJ-4LsRWC-1}\nwixd{NW1rTmOJ-3FfXTO-1}}}%
\nwixlogsorted{c}{{Serialize a term}{NW1rTmOJ-40h2U3-1}{\nwixu{NW1rTmOJ-4LsRWC-1}\nwixd{NW1rTmOJ-40h2U3-1}}}%
\nwixlogsorted{c}{{Sort.sml}{NW1rTmOJ-3ZqggC-1}{\nwixd{NW1rTmOJ-3ZqggC-1}}}%
\nwixlogsorted{c}{{Substitution of a term for a variable in a term}{NW1rTmOJ-3OIyM5-1}{\nwixu{NW1rTmOJ-4LsRWC-1}\nwixd{NW1rTmOJ-3OIyM5-1}}}%
\nwixlogsorted{c}{{Term.sml}{NW1rTmOJ-4LsRWC-1}{\nwixd{NW1rTmOJ-4LsRWC-1}}}%
\nwixlogsorted{c}{{Test if term contains subterm}{NW1rTmOJ-Re21r-1}{\nwixu{NW1rTmOJ-4LsRWC-1}\nwixd{NW1rTmOJ-Re21r-1}}}%
\nwixlogsorted{c}{{Test if two terms have the same sort}{NW1rTmOJ-1WQNps-1}{\nwixu{NW1rTmOJ-4LsRWC-1}\nwixd{NW1rTmOJ-1WQNps-1}}}%
\nwixlogsorted{i}{{\nwixident{Sort.CLASS}}{Sort.CLASS}}%
\nwixlogsorted{i}{{\nwixident{Sort.compare}}{Sort.compare}}%
\nwixlogsorted{i}{{\nwixident{Sort.FUN}}{Sort.FUN}}%
\nwixlogsorted{i}{{\nwixident{Sort.IND}}{Sort.IND}}%
\nwixlogsorted{i}{{\nwixident{Sort.t}}{Sort.t}}%
\nwixlogsorted{i}{{\nwixident{Sort.to{\_}string}}{Sort.to:unstring}}%
\nwixlogsorted{i}{{\nwixident{Term.compare}}{Term.compare}}%
\nwixlogsorted{i}{{\nwixident{Term.compare{\_}lists}}{Term.compare:unlists}}%
\nwixlogsorted{i}{{\nwixident{Term.eq}}{Term.eq}}%
\nwixlogsorted{i}{{\nwixident{Term.eq{\_}lists}}{Term.eq:unlists}}%
\nwixlogsorted{i}{{\nwixident{Term.fresh{\_}var}}{Term.fresh:unvar}}%
\nwixlogsorted{i}{{\nwixident{Term.fv}}{Term.fv}}%
\nwixlogsorted{i}{{\nwixident{Term.have{\_}same{\_}sort}}{Term.have:unsame:unsort}}%
\nwixlogsorted{i}{{\nwixident{Term.is{\_}class}}{Term.is:unclass}}%
\nwixlogsorted{i}{{\nwixident{Term.is{\_}fun}}{Term.is:unfun}}%
\nwixlogsorted{i}{{\nwixident{Term.is{\_}ind}}{Term.is:unind}}%
\nwixlogsorted{i}{{\nwixident{Term.occurs{\_}in}}{Term.occurs:unin}}%
\nwixlogsorted{i}{{\nwixident{Term.serialize}}{Term.serialize}}%
\nwixlogsorted{i}{{\nwixident{Term.sort}}{Term.sort}}%
\nwixlogsorted{i}{{\nwixident{Term.subst}}{Term.subst}}%
\nwixlogsorted{i}{{\nwixident{Term.t}}{Term.t}}%

